\section{Ba đối tượng của cùng một câu hỏi}
\label{sec:ch07-ba-doi-tuong}

\index{attribution!ba loại}%
\index{hàm attribution}%
Câu hỏi mở đầu chương cần một cách phát biểu đủ rộng để cả ba nhánh cùng nằm
trong. Zhang và cộng sự viết nó dưới dạng một \tn{hàm attribution}{attribution function}:
cho một mô hình $f$ đã huấn luyện và đầu ra $f(x)$ tại một đầu vào kiểm tra $x$,
hàm ấy gán một điểm số cho từng phần tử của một tập đối tượng đã
chọn~\autocite{p18unified}. Chọn tập nào chính là chỗ ba nhánh tách nhau; phần
còn lại của định nghĩa thì giống nhau.

\begin{definition}[Hàm attribution]
\label{def:ch07-ham-attribution}
Cho mô hình $f$, đầu vào kiểm tra $x$ và một tập đối tượng $E$. Một hàm
attribution $g$ gán cho mỗi phần tử $e \in E$ một số thực, được diễn giải là
mức đóng góp của $e$ vào đầu ra $f(x)$.
\end{definition}

Ba cách chọn $E$ dẫn tới ba nhánh. Khi $E$ là tập các đặc trưng của đầu vào, ta
có attribution đặc trưng, đối tượng của cả Phần~II và
chương~\ref{ch:attribution-tu-nguyen-ly-dau}, với điểm số $\varphi_j(x,f)$ cho
đặc trưng thứ $j$. Khi $E$ là tập dữ liệu huấn luyện
$\mathcal{D}_{\mathrm{train}}$, ta có \tn{attribution dữ liệu}{data attribution},
với điểm số $\psi_i(x)$ đo ảnh hưởng của điểm huấn luyện thứ $i$ lên đầu ra tại
$x$. Khi $E$ là tập các thành phần bên trong mô hình, chẳng hạn một nơ-ron, một
tầng hay một đầu chú ý, ta có \tn{attribution thành phần}{component attribution},
với điểm số $\gamma_k(x)$ cho thành phần $c_k$~\autocite{p18unified}.

Ba ký hiệu này giữ nguyên chữ cái mà bài báo dùng, nhưng chỉ số thì theo quy ước
của sách: $\varphi_j$ đã được chương~\ref{ch:attribution-tu-nguyen-ly-dau} gắn
với đặc trưng $j$, nên điểm dữ liệu nhận chỉ số $i$ và thành phần nhận chỉ số
$k$. Phụ lục~\ref{app:ky-hieu} ghi cả ba dòng cùng ký hiệu gốc trong bài báo.

Ba nhánh ấy phát triển trong ba cộng đồng khác nhau. Attribution đặc trưng thuộc
XAI, attribution dữ liệu thuộc nhánh lấy dữ liệu làm trung tâm, còn attribution
thành phần thuộc \tn{diễn giải cơ chế}{mechanistic interpretability}, hướng tìm
cách dựng lại cơ chế bên trong mô hình dưới dạng con người đọc
được~\autocite{p18unified}. Vì mỗi cộng đồng có hội nghị, khảo sát và thuật ngữ
riêng, một người đọc quen với LIME có thể không nhận ra rằng phương pháp
leave-one-out trong attribution dữ liệu đang dùng đúng \tn{phép nhiễu}{perturbation}
mà mình đã biết, chỉ khác là nó bỏ đi một điểm huấn luyện thay vì một đặc trưng.

Đây là một bài báo quan điểm, và các tác giả nói rõ như vậy: họ đề xuất một góc
nhìn, chứ không chứng minh một định lý hợp nhất~\autocite{p18unified}. Phần lớn
lập luận nằm ở chỗ chỉ ra cùng một kỹ thuật xuất hiện lại trong cả ba nhánh, nên
những kỹ thuật ấy là chỗ phải xem xét trước.
